\documentclass[10pt,aspectratio=169,serif]{beamer}

\usepackage[fontset=none, scheme=plain]{ctex}
\IfFontExistsTF{Songti SC}{%
	\setCJKmainfont{Songti SC}[AutoFakeSlant=0.2]%
	\setCJKsansfont{Songti SC}[AutoFakeSlant=0.2]%
}{%
	\setCJKmainfont{FandolSong-Regular}[Extension=.otf,BoldFont=FandolSong-Bold,AutoFakeSlant=0.2]%
	\setCJKsansfont{FandolSong-Regular}[Extension=.otf,BoldFont=FandolSong-Bold]%
}
\usepackage{amsmath,amssymb,graphicx,tikz,booktabs,multicol,calligra}
\usepackage{pgfplots}
\pgfplotsset{compat=1.18}
\usepackage{tcolorbox}
\tcbuselibrary{skins}
\usepackage{XyygU}
\usepackage{caption}
\captionsetup[figure]{name=图,labelsep=period}

\usetikzlibrary{arrows.meta,positioning,calc}

\setlength{\abovedisplayskip}{0pt plus 2pt minus 2pt}
\setlength{\belowdisplayskip}{0pt plus 2pt minus 2pt}

\title[5.2 三角函数的概念练习]{5.2 三角函数的概念（练习）}
\subtitle{三角函数定义 \& 同角三角函数基本关系}
\author{}
\institute{}
\date{2026-10-16}

\begin{document}

\begin{frame}
  \titlepage
\end{frame}

% ========== 5.2.1 练习 第1题 ==========
\begin{frame}[t]{5.2.1 练习第1题：填表}
\textbf{1.} 填表：
\vspace{0.5em}
\begin{center}
\begin{tabular}{|c|c|c|c|c|c|}
\hline
$\alpha$ & $2\pi$ & $\dfrac{13\pi}{6}$ & $-\pi$ & $-\dfrac{4\pi}{3}$ & $\dfrac{15\pi}{4}$ \\
\hline
$\sin\alpha$ & & & & & \\
\hline
$\cos\alpha$ & & & & & \\
\hline
$\tan\alpha$ & & & & & \\
\hline
\end{tabular}
\end{center}
\end{frame}

\begin{frame}[t]{5.2.1 练习第1题 解答}
\begin{tcolorbox}[colback=orange!10,colframe=orange,title=\textbf{解}]
先化到 $0\sim 2\pi$ 内终边相同的角：
\begin{itemize}
  \item $2\pi\equiv 0$：$\sin=0,\ \cos=1,\ \tan=0$；
  \item $\dfrac{13\pi}{6}=2\pi+\dfrac{\pi}{6}\equiv\dfrac{\pi}{6}$：$\sin=\dfrac{1}{2},\ \cos=\dfrac{\sqrt{3}}{2},\ \tan=\dfrac{\sqrt{3}}{3}$；
  \item $-\pi\equiv\pi$：$\sin=0,\ \cos=-1,\ \tan=0$；
  \item $-\dfrac{4\pi}{3}+2\pi=\dfrac{2\pi}{3}$：$\sin=\dfrac{\sqrt{3}}{2},\ \cos=-\dfrac{1}{2},\ \tan=-\sqrt{3}$；
  \item $\dfrac{15\pi}{4}-2\pi=\dfrac{7\pi}{4}$：$\sin=-\dfrac{\sqrt{2}}{2},\ \cos=\dfrac{\sqrt{2}}{2},\ \tan=-1$。
\end{itemize}
\end{tcolorbox}
\end{frame}

% ========== 5.2.1 练习 第2题 ==========
\begin{frame}[t]{5.2.1 练习第2题：三角形内角三角函数符号}
\textbf{2.}（口答）设 $\alpha$ 是三角形的一个内角，在 $\sin\alpha$，$\cos\alpha$，$\tan\alpha$，$\tan\dfrac{\alpha}{2}$ 中，哪些有可能取负值？
\end{frame}

\begin{frame}[t]{5.2.1 练习第2题 解答}
\begin{tcolorbox}[colback=orange!10,colframe=orange,title=\textbf{解}]
三角形内角 $0<\alpha<\pi$，即 $0^\circ<\alpha<180^\circ$。

\vspace{0.5em}
\begin{itemize}
  \item $\sin\alpha>0$（第一、二象限正弦均为正），不取负值；
  \item $\cos\alpha$：当 $90^\circ<\alpha<180^\circ$ 时为负，\alert{可能取负值}；
  \item $\tan\alpha$：当 $90^\circ<\alpha<180^\circ$ 时为负，\alert{可能取负值}；
  \item $0<\dfrac{\alpha}{2}<90^\circ$，$\tan\dfrac{\alpha}{2}>0$，不取负值。
\end{itemize}
\vspace{0.5em}
所以可能取负值的是 $\alert{\cos\alpha}$ 和 $\alert{\tan\alpha}$。
\end{tcolorbox}
\end{frame}

% ========== 5.2.1 练习 第3题 ==========
\begin{frame}[t]{5.2.1 练习第3题：确定三角函数值的符号}
\textbf{3.} 确定下列三角函数值的符号：
\begin{enumerate}
  \item $\sin 156^\circ$；
  \item $\cos\dfrac{16\pi}{5}$；
  \item $\cos(-450^\circ)$；
  \item $\tan\left(-\dfrac{17\pi}{8}\right)$；
  \item $\sin\left(-\dfrac{4\pi}{3}\right)$；
  \item $\tan 556^\circ$。
\end{enumerate}
\end{frame}

\begin{frame}[t]{5.2.1 练习第3题 解答}
\begin{tcolorbox}[colback=orange!10,colframe=orange,title=\textbf{解}]
\small
(1) $156^\circ$ 是第二象限角，$\sin 156^\circ \alert{>0}$。

(2) $\dfrac{16\pi}{5}=2\pi+\dfrac{6\pi}{5}$，$\dfrac{6\pi}{5}$ 是第三象限角，$\cos\dfrac{16\pi}{5}\alert{<0}$。

(3) $-450^\circ+2\times 360^\circ=-450^\circ+720^\circ=270^\circ$，终边在 $y$ 轴负半轴，$\cos(-450^\circ)=\cos 270^\circ=0$。

(4) $-\dfrac{17\pi}{8}+2\times 2\pi=-\dfrac{17\pi}{8}+\dfrac{32\pi}{8}=\dfrac{15\pi}{8}$，是第四象限角，$\tan\left(-\dfrac{17\pi}{8}\right)\alert{<0}$。

(5) $-\dfrac{4\pi}{3}+2\pi=\dfrac{2\pi}{3}$，是第二象限角，$\sin\left(-\dfrac{4\pi}{3}\right)\alert{>0}$。

(6) $556^\circ-360^\circ=196^\circ$，是第三象限角，$\tan 556^\circ\alert{>0}$。
\end{tcolorbox}
\end{frame}

% ========== 5.2.1 练习 第4题 ==========
\begin{frame}[t]{5.2.1 练习第4题：象限角充要条件}
\textbf{4.} 对于 ①$\sin\theta>0$，②$\sin\theta<0$，③$\cos\theta>0$，④$\cos\theta<0$，⑤$\tan\theta>0$ 与 ⑥$\tan\theta<0$，选择恰当的关系式序号填空：
\begin{enumerate}
  \item 角 $\theta$ 为第一象限角的充要条件是 \underline{\hspace{2cm}}；
  \item 角 $\theta$ 为第二象限角的充要条件是 \underline{\hspace{2cm}}；
  \item 角 $\theta$ 为第三象限角的充要条件是 \underline{\hspace{2cm}}；
  \item 角 $\theta$ 为第四象限角的充要条件是 \underline{\hspace{2cm}}。
\end{enumerate}
\end{frame}

\begin{frame}[t]{5.2.1 练习第4题 解答}
\begin{tcolorbox}[colback=orange!10,colframe=orange,title=\textbf{解}]
由各象限符号规律（一全正、二正弦、三正切、四余弦）：

\vspace{0.5em}
(1) 第一象限：$\sin>0$ 且 $\cos>0$（此时 $\tan>0$），填 \alert{①③}。

(2) 第二象限：$\sin>0$ 且 $\cos<0$（此时 $\tan<0$），填 \alert{①④}。

(3) 第三象限：$\sin<0$ 且 $\cos<0$（此时 $\tan>0$），填 \alert{②④}。

(4) 第四象限：$\sin<0$ 且 $\cos>0$（此时 $\tan<0$），填 \alert{②③}。
\end{tcolorbox}
\end{frame}

% ========== 5.2.1 练习 第5题 ==========
\begin{frame}[t]{5.2.1 练习第5题：求三角函数值}
\textbf{5.} 求下列三角函数值（可用计算工具，第(1)题精确到 $0.0001$）：
\begin{enumerate}
  \item $\cos 1109^\circ$；
  \item $\tan\dfrac{19\pi}{3}$；
  \item $\sin(-1050^\circ)$；
  \item $\tan\left(-\dfrac{31\pi}{4}\right)$。
\end{enumerate}
\end{frame}

\begin{frame}[t]{5.2.1 练习第5题 解答}
\begin{tcolorbox}[colback=orange!10,colframe=orange,title=\textbf{解}]
(1) $1109^\circ - 3\times 360^\circ = 29^\circ$，
$\cos 1109^\circ = \cos 29^\circ \approx \alert{0.8746}$。

(2) $\dfrac{19\pi}{3} - 4\times 2\pi = \dfrac{19\pi}{3}-\dfrac{24\pi}{3}=-\dfrac{5\pi}{3}\equiv\dfrac{\pi}{3}$，
$\tan\dfrac{19\pi}{3}=\tan\dfrac{\pi}{3}=\alert{\sqrt{3}}$。

(3) $-1050^\circ + 3\times 360^\circ = -1050^\circ+1080^\circ=30^\circ$，
$\sin(-1050^\circ)=\sin 30^\circ=\alert{\dfrac{1}{2}}$。

(4) $-\dfrac{31\pi}{4}+4\times 2\pi=-\dfrac{31\pi}{4}+\dfrac{32\pi}{4}=\dfrac{\pi}{4}$，
$\tan\left(-\dfrac{31\pi}{4}\right)=\tan\dfrac{\pi}{4}=\alert{1}$。
\end{tcolorbox}
\end{frame}

% ========== 5.2.2 练习 第1题 ==========
\begin{frame}[t]{5.2.2 练习第1题：已知余弦求正弦正切}
\textbf{1.} 已知 $\cos\alpha=-\dfrac{4}{5}$，且 $\alpha$ 为第三象限角，求 $\sin\alpha$，$\tan\alpha$ 的值。
\end{frame}

\begin{frame}[t]{5.2.2 练习第1题 解答}
\begin{tcolorbox}[colback=orange!10,colframe=orange,title=\textbf{解}]
由同角三角函数基本关系 $\sin^2\alpha+\cos^2\alpha=1$：
\[
\sin^2\alpha = 1-\cos^2\alpha = 1-\frac{16}{25} = \frac{9}{25}.
\]
因为 $\alpha$ 为第三象限角，$\sin\alpha<0$，所以
\[
\sin\alpha = -\frac{3}{5}.
\]
\[
\tan\alpha = \frac{\sin\alpha}{\cos\alpha} = \frac{-3/5}{-4/5} = \alert{\frac{3}{4}}.
\]
\end{tcolorbox}
\end{frame}

% ========== 5.2.2 练习 第2题 ==========
\begin{frame}[t]{5.2.2 练习第2题：已知正切求正弦余弦}
\textbf{2.} 已知 $\tan\varphi=-\sqrt{3}$，求 $\sin\varphi$，$\cos\varphi$ 的值。
\end{frame}

\begin{frame}[t]{5.2.2 练习第2题 解答}
\begin{tcolorbox}[colback=orange!10,colframe=orange,title=\textbf{解}]
由 $\tan\varphi=\dfrac{\sin\varphi}{\cos\varphi}=-\sqrt{3}$，得 $\sin\varphi=-\sqrt{3}\cos\varphi$。
代入 $\sin^2\varphi+\cos^2\varphi=1$：
\[
3\cos^2\varphi+\cos^2\varphi=1 \quad\Longrightarrow\quad \cos^2\varphi=\frac{1}{4}.
\]
$\tan\varphi<0$，$\varphi$ 在第二或第四象限，分两种情况：

\vspace{0.3em}
\textbf{第二象限：} $\cos\varphi=-\dfrac{1}{2}$，$\sin\varphi=\dfrac{\sqrt{3}}{2}$；

\textbf{第四象限：} $\cos\varphi=\dfrac{1}{2}$，$\sin\varphi=-\dfrac{\sqrt{3}}{2}$。
\end{tcolorbox}
\end{frame}

% ========== 5.2.2 练习 第3题 ==========
\begin{frame}[t]{5.2.2 练习第3题：已知正弦求余弦正切}
\textbf{3.} 已知 $\sin\theta=0.35$，求 $\cos\theta$，$\tan\theta$ 的值（精确到 $0.01$）。
\end{frame}

\begin{frame}[t]{5.2.2 练习第3题 解答}
\begin{tcolorbox}[colback=orange!10,colframe=orange,title=\textbf{解}]
\[
\cos^2\theta = 1-\sin^2\theta = 1-0.1225 = 0.8775,
\]
\[
\cos\theta = \pm\sqrt{0.8775}\approx\pm 0.94.
\]

\vspace{0.5em}
因为 $\sin\theta=0.35>0$，$\theta$ 在第一或第二象限：

\vspace{0.3em}
\textbf{第一象限：} $\cos\theta\approx 0.94$，$\tan\theta=\dfrac{0.35}{0.94}\approx 0.37$；

\textbf{第二象限：} $\cos\theta\approx -0.94$，$\tan\theta=\dfrac{0.35}{-0.94}\approx -0.37$。
\end{tcolorbox}
\end{frame}

% ========== 5.2.2 练习 第4题 ==========
\begin{frame}[t]{5.2.2 练习第4题：化简}
\textbf{4.} 化简：
\begin{enumerate}
  \item $\cos\theta\tan\theta$；
  \item $\dfrac{2\cos^2\alpha-1}{1-2\sin^2\alpha}$；
  \item $(1+\tan^2\alpha)\cos^2\alpha$。
\end{enumerate}
\end{frame}

\begin{frame}[t]{5.2.2 练习第4题 解答}
\begin{tcolorbox}[colback=orange!10,colframe=orange,title=\textbf{解}]
(1) $\cos\theta\tan\theta = \cos\theta\cdot\dfrac{\sin\theta}{\cos\theta} = \alert{\sin\theta}$。

(2) 由 $\cos^2\alpha+\sin^2\alpha=1$，
$2\cos^2\alpha-2\sin^2\alpha = 2(\cos^2\alpha-\sin^2\alpha)$，
分子 $2\cos^2\alpha-1=\cos^2\alpha-\sin^2\alpha$，
分母 $1-2\sin^2\alpha=\cos^2\alpha-\sin^2\alpha$，
所以原式 $=\alert{1}$。

(3) $(1+\tan^2\alpha)\cos^2\alpha = \cos^2\alpha+\sin^2\alpha = \alert{1}$。
\end{tcolorbox}
\end{frame}

% ========== 5.2.2 练习 第5题 ==========
\begin{frame}[t]{5.2.2 练习第5题：证明恒等式}
\textbf{5.} 求证：$\sin^4\alpha+\sin^2\alpha\cos^2\alpha+\cos^2\alpha=1$。
\end{frame}

\begin{frame}[t]{5.2.2 练习第5题 解答}
\begin{tcolorbox}[colback=orange!10,colframe=orange,title=\textbf{证明}]
左边提取公因式：
\[
\sin^4\alpha+\sin^2\alpha\cos^2\alpha+\cos^2\alpha
= \sin^2\alpha(\sin^2\alpha+\cos^2\alpha)+\cos^2\alpha.
\]
由 $\sin^2\alpha+\cos^2\alpha=1$：
\[
= \sin^2\alpha\cdot 1 + \cos^2\alpha = \sin^2\alpha+\cos^2\alpha = 1 = \text{右边}.
\]
\vspace{0.3em}
故原式成立。
\end{tcolorbox}
\end{frame}

\begin{frame}[c]
  \begin{center}
    {\Huge\calligra Thanks!}
  \end{center}
\end{frame}

\end{document}
